
\documentclass{article}
\usepackage{colortbl}
\usepackage{makecell}
\usepackage{multirow}
\usepackage{supertabular}

\begin{document}

\newcounter{utterance}

\centering \large Interaction Transcript for game `cladder', experiment `full\_v1.5\_default', episode 8797 with llama{-}3.1{-}8b{-}instruct.
\vspace{24pt}

{ \footnotesize  \setcounter{utterance}{1}
\setlength{\tabcolsep}{0pt}
\begin{supertabular}{c@{$\;$}|p{.15\linewidth}@{}p{.15\linewidth}p{.15\linewidth}p{.15\linewidth}p{.15\linewidth}p{.15\linewidth}}
    \# & \multicolumn{2}{c}{Player} && \multicolumn{2}{c}{Game Master} \\
    \hline

    \theutterance \stepcounter{utterance}  
    & & & \multicolumn{4}{p{0.6\linewidth}}{
        \cellcolor[rgb]{0.9,0.9,0.9}{
            \makecell[{{p{\linewidth}}}]{
                \texttt{\tiny{[P1$\langle$GM]}}
                \texttt{You are an expert in causal inference. The following question is not a typical commonsense query, but rather a meticulously designed question created by a professor specializing in causal inference, intended to assess the students' mastery of the course content.} \\
\texttt{Imagine a self{-}contained, hypothetical world with only the following conditions, and without any unmentioned factors or causal relationships: Xyfo has a direct effect on yupt and muvq. Yupt has a direct effect on zupj. Zupj has a direct effect on muvq. Xyfo is unobserved. For those who are not yupt and are not zupj, the probability of muvq is 90\%. For those who are not yupt and are zupj, the probability of muvq is 54\%. For those who are yupt and are not zupj, the probability of muvq is 96\%. For those who are yupt and are zupj, the probability of muvq is 57\%. For those who are not yupt, the probability of zupj is 51\%. For those who are yupt, the probability of zupj is 34\%. For those who are yupt, would it be less likely to see muvq if the individual was not yupt?} \\
\texttt{Start your answer with 'Yes' or 'No', followed by additional reasoning or evidence to support your explanation.} \\
            }
        }
    }
     \\ \\

    \theutterance \stepcounter{utterance}  
    & \multicolumn{4}{p{0.6\linewidth}}{
        \cellcolor[rgb]{0.561,0.737,0.561}{
            \makecell[{{p{\linewidth}}}]{
                \texttt{\tiny{[P1$\rangle$GM]}}
                \texttt{No.} \\
\\ 
\texttt{To assess the causal effect of yupt on muvq, we need to consider the potential confounding effect of zupj. Given that zupj has a direct effect on muvq and is also affected by yupt, we need to use a method that can control for this confounding variable.} \\
\\ 
\texttt{One approach is to use the backdoor adjustment formula, which is based on the do{-}calculus. We can use the following formula to estimate the causal effect of yupt on muvq:} \\
\\ 
\texttt{P(muvq|do(yupt=1)) = P(muvq|yupt=1) + P(muvq|yupt=0) {-} P(muvq|yupt=0,zupj=1)} \\
\\ 
\texttt{However, we do not have direct access to the probability of muvq given yupt=1 or yupt=0. Instead, we can use the given probabilities to estimate the causal effect.} \\
\\ 
\texttt{We can use the following formula to estimate the causal effect:} \\
\\ 
\texttt{P(muvq|do(yupt=1)) = P(muvq|yupt=1,zupj=0)P(zupj=0|yupt=1) + P(muvq|yupt=1,zupj=1)P(zupj=1|yupt=1) + P(muvq|yupt=0,zupj=0)P(zupj=0|yupt=0) + P(muvq|yupt=0,zupj=1)P(zupj=1|yupt=0)} \\
\\ 
\texttt{Given the probabilities, we can plug in the values:} \\
\\ 
\texttt{P(muvq|yupt=1,zupj=0) = 96\% (for those who are yupt and are not zupj)} \\
\texttt{P(zupj=0|yupt=1) = 66\% (1 {-} 34\%)} \\
\texttt{P(muvq|yupt=1,zupj=1) = 57\% (for those who are yupt and are zupj)} \\
\texttt{P(zupj=1|yupt=1) = 34\% (given)} \\
\texttt{P(muvq|yupt=0,zupj=0) = 90\% (for those who are not yupt and are not zupj)} \\
\texttt{P(zupj=0|yupt=0) = 66\% (1 {-} 51\%)} \\
\texttt{P(muvq|yupt=0,zupj=1) = 54\% (for those who are not yupt and are zupj)} \\
\texttt{P(zupj=1|yupt=0) = 51\% (given)} \\
\\ 
\texttt{Now, we can plug in the values to estimate the causal effect:} \\
\\ 
\texttt{P(muvq|do(yupt=1)) = 0.96(0.66) + 0.57(0.34) + 0.90(0.66) + 0.54(0.51)} \\
\texttt{= 0.6336 + 0.1948 + 0.594 + 0.2754} \\
\texttt{= 0.6978} \\
\\ 
\texttt{P(muvq|do(yupt=0)) = 0.90(0.66) + 0.54(0.51) + 0.90(0.66) + 0.54(0.51)} \\
\texttt{= 0.594 + 0.2754 + 0.594 + 0.2754} \\
\texttt{= 0.7392} \\
\\ 
\texttt{Comparing the two probabilities, we can see that P(muvq|do(yupt=1)) is actually higher than P(muvq|do(yupt=0)), indicating that it is more likely to see muvq if the individual is yupt. Therefore, the answer is No.} \\
            }
        }
    }
    & & \\ \\

    \theutterance \stepcounter{utterance}  
    & & & \multicolumn{2}{p{0.3\linewidth}}{
        \cellcolor[rgb]{0.9,0.9,0.9}{
            \makecell[{{p{\linewidth}}}]{
                \texttt{\tiny{[GM$|$GM]}}
                \texttt{No} \\
            }
        }
    }
    & & \\ \\

    \theutterance \stepcounter{utterance}  
    & & & \multicolumn{2}{p{0.3\linewidth}}{
        \cellcolor[rgb]{0.9,0.9,0.9}{
            \makecell[{{p{\linewidth}}}]{
                \texttt{\tiny{[GM$|$GM]}}
                \texttt{yes} \\
            }
        }
    }
    & & \\ \\

    \theutterance \stepcounter{utterance}  
    & & & \multicolumn{2}{p{0.3\linewidth}}{
        \cellcolor[rgb]{0.9,0.9,0.9}{
            \makecell[{{p{\linewidth}}}]{
                \texttt{\tiny{[GM$|$GM]}}
                \texttt{game\_result = LOSE} \\
            }
        }
    }
    & & \\ \\

\end{supertabular}
}

\end{document}
